\documentclass[12pt,a4paper]{article}

\usepackage[a4paper,margin=1in]{geometry}
\usepackage{graphicx}
\usepackage{booktabs}
\usepackage{amsmath,amssymb}
\usepackage{float}
\usepackage{hyperref}
\usepackage{xcolor}
\usepackage{listings}
\usepackage{enumitem}
\usepackage{fancyhdr}
\usepackage{longtable}

\hypersetup{colorlinks=true,linkcolor=blue,urlcolor=blue}

\pagestyle{fancy}
\fancyhf{}
\lhead{ICS1512}
\rhead{Machine Learning Algorithms Laboratory}
\cfoot{\thepage}

\lstset{
language=Python,
basicstyle=\ttfamily\small,
numbers=left,
frame=single,
breaklines=true,
showstringspaces=false,
keywordstyle=\color{blue},
commentstyle=\color{gray},
stringstyle=\color{purple}
}

\begin{document}

\begin{tabular}{|p{4cm}|p{11cm}|}
\hline
Experiment No. & 5 \\ \hline
Experiment Title & Decision Tree and Random Forest: A Comparative Classification Study
\footnote{\textbf{GitHub Repository: }\url{https://github.com/Nightingales21/ICS1512-Machine-Learning-Labratory}}\\ \hline
Student Name & Nitin \\ \hline
Register Number & 312224700142 \\ \hline
Faculty & Dr. Poreddy Ajay Kumar Reddy \\ \hline
Submission Date & \today \\ \hline
\end{tabular}

\section{Objective}
\begin{itemize}
\item To implement a Decision Tree classifier.
\item To extend the Decision Tree into a Random Forest ensemble model.
\item To study the impact of hyperparameters on overfitting and generalization.
\item To select optimal hyperparameters using 5-Fold Cross-Validation.
\item To compare single-tree and ensemble-tree models.
\end{itemize}

\section{Dataset Description}
\begin{longtable}{|p{6cm}|p{8cm}|}
\hline
Dataset Name & Wisconsin Diagnostic Breast Cancer (WDBC) \\ \hline
Dataset Source & UCI Machine Learning Repository (loaded via
\texttt{sklearn.datasets.load\_breast\_cancer}, an identical copy of the
canonical UCI WDBC data) \\ \hline
Number of Samples & 569 \\ \hline
Number of Features & 30 numerical attributes (mean, standard-error, and
``worst'' value of 10 base measurements per cell nucleus) \\ \hline
Number of Classes & 2 (Malignant, Benign) \\ \hline
Missing Values & 0 (verified programmatically) \\ \hline
Train--Test Split & 80\% / 20\%, stratified, \texttt{random\_state=42} \\ \hline
\end{longtable}

Class balance: 357 benign samples (62.7\%) and 212 malignant samples
(37.3\%) -- a moderate imbalance. \textbf{Label encoding note:} scikit-learn's
built-in loader encodes \texttt{target = 0} as \emph{malignant} and
\texttt{target = 1} as \emph{benign}; this (reversed from the
0=negative/1=positive convention some readers expect) is used consistently
throughout this report and all confusion matrices below.

\section{Preprocessing Steps}
\begin{itemize}
\item \textbf{Class label encoding:} The dataset's numeric \texttt{target}
column (0 = malignant, 1 = benign) was mapped to a readable
\texttt{diagnosis} column for EDA and reporting purposes; the numeric
encoding was used directly for model training.
\item \textbf{Missing value handling:} None required; the dataset was
verified to contain zero missing cells.
\item \textbf{Feature scaling:} Not required. Both Decision Tree and Random
Forest are threshold-split-based models and are invariant to monotonic
feature scaling, so no \texttt{StandardScaler} step was applied (unlike the
distance-/margin-based models in Experiments 2 and 4).
\item \textbf{Train--test split:} \texttt{train\_test\_split} with
\texttt{test\_size=0.2}, \texttt{stratify=y}, \texttt{random\_state=42}.
\end{itemize}

\section{Exploratory Data Analysis}

The reusable function \texttt{generate\_eda\_summary()} from
\texttt{ml\_lab\_utils.py} (Experiment 1, unmodified) was used to produce
the consolidated 12-panel EDA figure below.

\begin{figure}[H]
\centering
\includegraphics[width=\linewidth]{figures/eda_breast_cancer.png}
\caption{Consolidated 12-panel EDA summary for the Wisconsin Breast Cancer dataset.}
\label{fig:eda}
\end{figure}

\textbf{Inference:} The dataset has no missing values and a moderate class
imbalance (62.7\% benign, 37.3\% malignant). \texttt{worst area} and
\texttt{mean area} are the two highest-variance features (Panels 6, 9, 10,
12) and are, unsurprisingly, very strongly correlated with each other
(Panel 9) since both derive from cell/nucleus size. The box plot (Panel 7)
shows several high-value outliers among the size-related features, and the
Q--Q plot (Panel 10) confirms \texttt{worst area} is right-skewed rather
than normally distributed -- consistent with a small subset of unusually
large, aggressive tumors in the dataset.

\begin{figure}[H]
\centering
\includegraphics[width=\linewidth]{figures/full_correlation_heatmap.png}
\caption{Full 30-feature correlation heatmap, required per Section 2 of the lab manual (Feature correlation analysis).}
\label{fig:full_corr}
\end{figure}
\textbf{Inference:} The heatmap reveals strong block-diagonal correlation
structure: the \texttt{mean}, \texttt{error}, and \texttt{worst} versions of
\texttt{radius}, \texttt{perimeter}, and \texttt{area} are all very highly
correlated with each other (dark red blocks), as expected since perimeter
and area are geometrically derived from radius. Similarly,
\texttt{concavity}, \texttt{concave points}, and \texttt{compactness} form a
second correlated cluster. This multicollinearity means many of the 30
raw features are redundant, which explains why both models below achieve
strong performance using only a handful of the most informative features
(Section~10, Feature Importance).

\section{Implementation Details}

\textbf{Decision Tree} recursively partitions the feature space by choosing,
at each node, the feature and threshold that most reduces class impurity
(measured by \textbf{Gini index} or \textbf{entropy}). Deeper trees can fit
the training data almost perfectly but are prone to high variance and
overfitting, since individual splits become sensitive to noise in small
leaf subsets.

\textbf{Random Forest} is an ensemble of many decision trees, each trained
on a bootstrap-resampled subset of the training data (\textbf{bagging}) and
restricted to a random subset of features at each split
(\textbf{random feature selection}). Final predictions are made by
\textbf{majority vote} across all trees. Averaging over many
de-correlated trees reduces the variance of the ensemble relative to any
single tree, generally improving generalization at the cost of reduced
interpretability and increased training/prediction cost.

\begin{lstlisting}[language=Python, caption={Baseline Decision Tree via the reusable classification function, from \texttt{ml\_lab\_utils.py} (Experiment 1).}]
baseline_models = {"Decision Tree (Baseline)": DecisionTreeClassifier(random_state=42)}
baseline_results_df, baseline_fitted = train_evaluate_classification(
    baseline_models, X_train, X_test, y_train, y_test, scale=False)
\end{lstlisting}

\begin{lstlisting}[language=Python, caption={Hyperparameter search spaces and 5-fold GridSearchCV, from \texttt{experiment5\_dt\_rf\_breast\_cancer.py}.}]
dt_param_grid = {
    "criterion": ["gini", "entropy"],
    "max_depth": [3, 5, 7, 10, None],
    "min_samples_split": [2, 5, 10],
    "min_samples_leaf": [1, 2, 4],
}
rf_param_grid = {
    "n_estimators": [50, 100, 200],
    "max_depth": [5, 10, None],
    "max_features": ["sqrt", "log2"],
    "bootstrap": [True, False],
}
dt_grid = GridSearchCV(DecisionTreeClassifier(random_state=42), dt_param_grid,
                        cv=StratifiedKFold(5, shuffle=True, random_state=42),
                        scoring="accuracy")
\end{lstlisting}

\section{Experimental Setup}
\begin{tabular}{ll}
\toprule
Python Version & 3.12 \\
Scikit-Learn Version & 1.8.0 \\
IDE & Jupyter Notebook / VS Code \\
Operating System & Ubuntu 24.04 LTS \\
Hardware & Standard lab workstation (CPU only) \\
\bottomrule
\end{tabular}

\vspace{0.3cm}
\textbf{Environment activation command (as mandated by the lab manual):}
\begin{lstlisting}[language=bash]
source /opt/anaconda3/bin/activate anaconda-navigation
\end{lstlisting}

\section{Hyperparameter Tuning Results}

\textbf{Table 1: Decision Tree Hyperparameter Evaluation using 5-Fold Cross-Validation}
\begin{tabular}{lccc}
\toprule
Criterion & Max Depth & Avg CV Accuracy (\%) & Avg CV F1 Score \\
\midrule
gini    & 3    & 93.19 & 0.9455 \\
gini    & 5    & \textbf{93.85} & \textbf{0.9506} \\
gini    & 7    & 93.85 & 0.9506 \\
gini    & 10   & 93.85 & 0.9506 \\
gini    & None & 93.85 & 0.9506 \\
entropy & 3    & 93.41 & 0.9480 \\
entropy & 5    & 92.97 & 0.9439 \\
entropy & 7    & 92.97 & 0.9435 \\
entropy & 10   & 92.97 & 0.9435 \\
entropy & None & 92.97 & 0.9435 \\
\bottomrule
\end{tabular}

\vspace{0.2cm}
\noindent \textbf{Overall best Decision Tree (GridSearchCV, full 4-parameter
search):} criterion=gini, max\_depth=5, min\_samples\_leaf=4,
min\_samples\_split=2, CV accuracy = 93.85\%. Note that accuracy plateaus
identically from max\_depth=5 onward within the gini rows above because the
\texttt{min\_samples\_leaf=4} constraint (selected jointly by the full grid
search) naturally stops the tree from growing much beyond depth 5 regardless
of the nominal \texttt{max\_depth} ceiling.

\vspace{0.3cm}
\textbf{Table 2: Random Forest Hyperparameter Evaluation using 5-Fold Cross-Validation}
\begin{tabular}{lcccc}
\toprule
n\_estimators & Max Depth & Max Features & Avg CV Accuracy (\%) & Avg CV F1 \\
\midrule
50  & 5    & sqrt & 96.04 & 0.9687 \\
50  & 10   & log2 & 96.26 & 0.9702 \\
50  & None & log2 & 96.26 & 0.9702 \\
100 & 5    & sqrt & 96.04 & 0.9685 \\
100 & 10   & log2 & \textbf{96.48} & \textbf{0.9718} \\
100 & None & log2 & 96.48 & 0.9718 \\
200 & 5    & log2 & 95.82 & 0.9666 \\
200 & 10   & sqrt & 96.26 & 0.9700 \\
200 & None & sqrt & 96.26 & 0.9700 \\
\bottomrule
\end{tabular}

\vspace{0.2cm}
\noindent \textbf{Overall best Random Forest (GridSearchCV, full 4-parameter
search):} bootstrap=True, max\_depth=10, max\_features=log2,
n\_estimators=100, CV accuracy = 96.48\%.

\vspace{0.3cm}
\textbf{Hyperparameter Tuning Results Summary}
\begin{tabular}{lllc}
\toprule
Model & Search Method & Best Parameters & Best CV Accuracy \\
\midrule
Decision Tree & GridSearchCV (5-fold) & gini, depth=5, leaf=4, split=2 & 0.9385 \\
Random Forest & GridSearchCV (5-fold) & bootstrap=T, depth=10, log2, n=100 & \textbf{0.9648} \\
\bottomrule
\end{tabular}

\section{Decision Tree Performance}
\begin{tabular}{lc}
\toprule
Metric & Value \\
\midrule
Accuracy          & 0.9035 \\
Precision         & 0.9061 \\
Recall            & 0.9035 \\
F1 Score          & 0.9041 \\
ROC-AUC           & 0.9358 \\
Training Time (s) & 0.0059 \\
\bottomrule
\end{tabular}

\begin{figure}[H]
\centering
\includegraphics[width=0.7\linewidth]{figures/cm_decision_tree.png}
\caption{Confusion matrix for the tuned Decision Tree on the held-out test set.}
\label{fig:cm_dt}
\end{figure}

\begin{figure}[H]
\centering
\includegraphics[width=\linewidth]{figures/decision_tree_structure.png}
\caption{Structure of the tuned Decision Tree (first 3 levels shown for readability).}
\label{fig:tree_structure}
\end{figure}
\textbf{Inference:} The tree's root split is on \texttt{worst radius}
($\leq 16.795$), immediately separating a large, highly-pure benign-leaning
branch (277 of 304 benign) from a malignant-leaning branch (143 of 151
malignant). This single split alone achieves strong separation, confirming
that tumor size (captured by \texttt{worst radius}) is the single most
discriminative measurement in this dataset -- consistent with the feature
importance ranking in Section~10.

\section{Overfitting Analysis: Tree Depth vs. Accuracy}
\begin{figure}[H]
\centering
\includegraphics[width=0.8\linewidth]{figures/depth_vs_accuracy.png}
\caption{Training vs. 5-fold validation accuracy as a function of maximum tree depth.}
\label{fig:depth}
\end{figure}
\textbf{Inference:} Training accuracy climbs monotonically to a perfect
1.000 by depth 10 and beyond, while validation accuracy peaks early, around
depth 2 (0.932), and then \emph{declines} as depth increases further,
falling to 0.916 by depth 10 and staying there for all deeper trees. This
divergence -- training accuracy rising while validation accuracy falls -- is
the textbook signature of overfitting: beyond a depth of roughly 2--4, the
tree begins memorizing noise specific to the training folds rather than
learning generalizable decision rules.

\section{Random Forest Performance}
\begin{tabular}{lc}
\toprule
Metric & Value \\
\midrule
Accuracy          & 0.9561 \\
Precision         & 0.9561 \\
Recall            & 0.9561 \\
F1 Score          & 0.9560 \\
ROC-AUC           & 0.9924 \\
Training Time (s) & 0.1426 \\
\bottomrule
\end{tabular}

\begin{figure}[H]
\centering
\includegraphics[width=0.7\linewidth]{figures/cm_random_forest.png}
\caption{Confusion matrix for the tuned Random Forest on the held-out test set.}
\label{fig:cm_rf}
\end{figure}

\begin{figure}[H]
\centering
\includegraphics[width=0.7\linewidth]{figures/n_estimators_vs_accuracy.png}
\caption{Random Forest: 5-fold CV accuracy as a function of ensemble size (n\_estimators).}
\label{fig:n_est}
\end{figure}
\textbf{Inference:} Accuracy improves quickly from a single tree (0.925,
identical to an ordinary Decision Tree) up to about 10--50 trees (0.956),
then continues to improve more slowly, plateauing around 150--300 trees
(0.965). This is the expected bagging behavior: variance reduction from
averaging has diminishing but still-positive returns as more
(de-correlated) trees are added, with most of the benefit realized within
the first 50--100 trees.

\section{5-Fold Cross-Validation Performance Comparison}

\textbf{Table 3: 5-Fold Cross-Validation Accuracy Comparison}
\begin{tabular}{lccccccc}
\toprule
Model & Fold 1 & Fold 2 & Fold 3 & Fold 4 & Fold 5 & Average \\
\midrule
Decision Tree  & 0.9451 & 0.9451 & 0.9121 & 0.9231 & 0.9670 & 0.9385 \\
Random Forest  & 0.9670 & 0.9780 & 0.9341 & 0.9560 & 0.9890 & \textbf{0.9648} \\
\bottomrule
\end{tabular}

\begin{figure}[H]
\centering
\includegraphics[width=0.75\linewidth]{figures/cv_accuracy_comparison.png}
\caption{5-fold cross-validation accuracy: Decision Tree vs. Random Forest.}
\label{fig:cv}
\end{figure}
\textbf{Inference:} Random Forest outperforms the Decision Tree on
\emph{every single fold}, not just on average -- the two curves never
cross. This fold-by-fold consistency is stronger evidence of a genuine
improvement than the average alone, since it rules out the possibility that
Random Forest's advantage is driven by a single easy fold.

\section{Evaluation Metrics: ROC Curve and Comparative Analysis}
\begin{figure}[H]
\centering
\includegraphics[width=\linewidth]{figures/roc_curves.png}
\caption{ROC curves for tuned Decision Tree and tuned Random Forest.}
\label{fig:roc}
\end{figure}
\textbf{Inference:} Random Forest's ROC curve dominates the Decision Tree's
at every false-positive rate, with a substantially higher AUC (0.9924 vs.\
0.9358). For a medical diagnosis task such as this one, where both false
negatives (missing a malignant tumor) and false positives (unnecessary
biopsy/anxiety for a benign case) carry real costs, this gap represents a
practically meaningful improvement in the model's ability to rank malignant
cases above benign ones across all possible decision thresholds.

\begin{figure}[H]
\centering
\includegraphics[width=\linewidth]{figures/model_comparison_bar.png}
\caption{Decision Tree vs. Random Forest: Accuracy, Precision, Recall, and F1 comparison.}
\label{fig:bar}
\end{figure}

\begin{figure}[H]
\centering
\includegraphics[width=0.8\linewidth]{figures/feature_importance_comparison.png}
\caption{Feature importance comparison: Decision Tree vs. Random Forest (top 10 features).}
\label{fig:importance}
\end{figure}
\textbf{Inference:} The single Decision Tree concentrates almost all of its
importance (0.75) on just one feature, \texttt{worst radius}, using only 5
features with non-zero importance in total -- a direct consequence of
greedy, single-path splitting. Random Forest spreads importance far more
evenly across 10+ features (\texttt{worst concave points},
\texttt{worst area}, \texttt{worst radius}, \texttt{mean perimeter}, etc.,
each contributing roughly 5--13\%), because each tree in the ensemble is
restricted to a random feature subset at every split, forcing the forest as
a whole to discover and rely on multiple complementary, correlated but
non-identical signals rather than collapsing onto a single dominant
feature.

\vspace{0.3cm}
\begin{tabular}{lcc}
\toprule
Criterion & Decision Tree & Random Forest \\
\midrule
Accuracy (test split) & 0.9035 & \textbf{0.9561} \\
Accuracy (5-fold CV)   & 0.9385 & \textbf{0.9648} \\
ROC-AUC                & 0.9358 & \textbf{0.9924} \\
Model Complexity       & Low & High \\
Training Time          & Low (0.006\,s) & Low--Moderate (0.143\,s) \\
Interpretability       & High & Low \\
\bottomrule
\end{tabular}

\section{Observations}
\begin{itemize}
\item \textbf{How does tree depth affect overfitting?} As shown in
Figure~\ref{fig:depth}, increasing depth beyond about 2--4 causes training
accuracy to keep rising toward a perfect 1.0 while validation accuracy
peaks and then declines -- the tree is fitting noise specific to individual
training folds (high variance), a textbook overfitting pattern.
\item \textbf{Which hyperparameter had the greatest impact on
performance?} For the Decision Tree, \texttt{max\_depth} (in combination
with \texttt{min\_samples\_leaf}) had the largest effect, moving CV
accuracy by roughly 1.5 percentage points across the tested range and
directly controlling the overfitting behavior above; \texttt{criterion}
(gini vs.\ entropy) had a much smaller effect (under 1 percentage point).
For Random Forest, \texttt{n\_estimators} had the most visible effect at
the low end (Figure~\ref{fig:n_est}: +3 points from 1 to 10 trees), while
\texttt{max\_depth}, \texttt{max\_features}, and \texttt{bootstrap} produced
smaller, second-order refinements once the ensemble was reasonably large.
\item \textbf{How does Random Forest improve generalization?} By
averaging predictions across many trees, each trained on a different
bootstrap sample and restricted to a random feature subset per split,
Random Forest decorrelates the individual trees' errors. Errors that are
independent across trees partially cancel out under majority voting,
reducing the ensemble's overall variance without a corresponding increase
in bias -- visible here as Random Forest beating the Decision Tree on
\emph{every} cross-validation fold (Table 3) despite each individual tree
in the forest being no more sophisticated than the single tuned Decision
Tree.
\item \textbf{Did ensemble learning always improve performance? Why or why
not?} In this experiment, yes -- Random Forest improved every metric
tracked (accuracy, precision, recall, F1, ROC-AUC) and every cross-
validation fold over the Decision Tree. This is not universally guaranteed:
ensemble gains are largest when the base learners are prone to high
variance (as deep, unconstrained decision trees are) and when there is
enough data to support many bootstrap resamples; on datasets with very few
samples, strong label noise, or where a single well-regularized tree
already captures the true decision boundary well, the added complexity and
reduced interpretability of a forest may not be worth a marginal (or
negligible) accuracy gain.
\end{itemize}

\section{Conclusion}
Decision Tree and Random Forest models were implemented and evaluated on
the Wisconsin Diagnostic Breast Cancer dataset using 5-fold cross-
validation for hyperparameter selection. The tuned Decision Tree
(gini, max\_depth=5, min\_samples\_leaf=4) achieved 93.85\% CV accuracy and
90.35\% test accuracy, while the tuned Random Forest (100 trees, max\_depth=10,
max\_features=log2, bootstrapping enabled) achieved 96.48\% CV accuracy and
95.61\% test accuracy -- a consistent, fold-by-fold improvement across every
evaluation metric, most notably raising ROC-AUC from 0.936 to 0.992. These
results directly confirm the theoretical expectation that bagging and
random feature selection reduce variance and improve generalization
relative to a single decision tree, at the cost of reduced interpretability
(a single readable tree vs.\ an opaque 100-tree ensemble) and a roughly
24$\times$ increase in training time (still well under a fifth of a second
on this dataset size, so practically negligible here). For this diagnostic
task, where predictive reliability is paramount, Random Forest is the
clearly preferable model; the single Decision Tree remains valuable
primarily as an interpretable baseline and as a direct visualization of the
dataset's dominant decision rule (Figure~\ref{fig:tree_structure}).

\section{References}
\begin{enumerate}
\item Scikit-learn Developers, ``Decision Trees,'' scikit-learn
documentation. [Online]. Available:
\url{https://scikit-learn.org/stable/modules/tree.html}
\item Scikit-learn Developers, ``Ensemble methods: Random Forests,''
scikit-learn documentation. [Online]. Available:
\url{https://scikit-learn.org/stable/modules/ensemble.html#random-forests}
\item W. H. Wolberg, W. N. Street, and O. L. Mangasarian, ``Breast Cancer
Wisconsin (Diagnostic) Data Set,'' UCI Machine Learning Repository, 1995.
[Online]. Available: \url{https://archive.ics.uci.edu}
\end{enumerate}

\end{document}
